% ----------------------------------------------------------------
% ----------------------------------------------------------------
% AMS-LaTeX Paper ************************************************f
% **** -----------------------------------------------------------
% documentclass[fleqn]{report}
%\RequirePackage{silence}
%\WarningFilter{lualatex}{Token not allowed}
% Einzelne Warnung eines Pakets unterdrücken:
%\WarningFilter{pdftex}{destination with}

% Alle Warnungen eines bestimmten Pakets unterdrücken:
%\WarningsOff[lualatex]
\PassOptionsToPackage{unicode}{hyperref}
%\PassOptionsToPackage{naturalnames}{hyperref}
\documentclass[a4paper,fleqn]{article}
%\documentclass[a4paper,fleqn]{scrreprt}
%\documentclass{styles/wissdoc-kw-eng}
\setlength{\oddsidemargin}{0.0cm}
\setlength{\evensidemargin}{0.0cm}
\setlength{\paperheight}{297mm}
%\usepackage[utf8]{inputenc}
\usepackage{german}
\usepackage{graphicx} 
\usepackage{amsmath}
\usepackage{amssymb}
\usepackage{alltt}
\usepackage{color}
\usepackage{fancyvrb}
\usepackage{verbatim}
\usepackage{mathtools}
\DeclarePairedDelimiter\bra{\langle}{\rvert}
\DeclarePairedDelimiter\ket{\lvert}{\rangle}
\DeclarePairedDelimiterX\braket[2]{\langle}{\rangle}{#1\,\delimsize\vert\,\mathopen{}#2}
\DeclarePairedDelimiterX\expect[1]{\langle}{\rangle}{#1}
%\newcounter{thm}[section]
%\newtheorem{theorem}[thm]{thm}
%\newtheorem{lemma}[theorem]{Lemma}
%\newtheorem{proposition}[theorem]{Hilfssatz}
%\newtheorem{corollary}[theorem]{Corollary}
%\newtheorem{definition}[theorem]{Definition}
%\newenvironment{theorem}[Theorem]{\begin{trivlist}
%\item[\hskip \labelsep {\textit \textit{#1}}]}{\end{trivlist}}
%\newenvironment{proof}[1][Beweis]{\begin{trivlist}
%\item[\hskip \labelsep {\bfseries #1}]}{\end{trivlist}}
%\item[\hskip \labelsep {\bfseries #1}]}{\end{trivlist}}
\newcommand*\codeRed[1]{\textcolor[rgb]{1,0,0}{\textbf{#1}}}
\newcommand*\codeGreen[1]{\textcolor[rgb]{0,0.5,0}{\textbf{#1}}}
\newcommand{\lb}{[}
\newcommand{\rb}{]}
\newcommand{\codeCall}[1]{\textbf{\texttt{#1}}}
\newcommand{\mcL}{\mathcal{L}}
\newcommand{\kb}{k_{\mathrm{B}}}
\newcommand{\titlemath}[1]{\texorpdfstring{#1}{}}
\newcommand{\pmat}[1]{\begin{pmatrix}#1\end{pmatrix}}


%\newenvironment{\vspace{1ex}\noindent{\bf Proof}\hspace{0.5em}}
%	{\hfill\qed\vspace{1ex}}
 

\newenvironment{example}[1][Example]{\begin{trivlist}
\item[\hskip \labelsep {\bfseries #1}]}{\end{trivlist}}
\newenvironment{remark}[1][Remark]{\begin{trivlist}
\item[\hskip \labelsep {\bfseries #1}]}{\end{trivlist}}

%\newcommand{\qed}{\nobreak \ifvmode \relax \else
%      \ifdim\lastskip<1.5em \hskip-\lastskip
%      \hskip1.5em plus0em minus0.5em \fi \nobreak
%      \vrule height0.75em width0.5em depth0.25em\fi}
%\newcommand{\qed}{$\square$}
\newcommand{\qt}{\texttt}

\newcommand{\bibdef}[4]{\bibitem{#1} {\bf #2:} {\it #3}\\#4}


\newcommand{\gauss}[1]{\lfloor #1\rfloor}
\numberwithin{equation}{section}


%\setcounter{equation}{1}

\setcounter{tocdepth}{4}
\setcounter{secnumdepth}{4}

\newcommand{\ltri}{\vartriangleleft}
\usepackage{amsthm}
\newtheorem{theorem}{Theorem}[section]
\renewcommand{\proofname}{Beweis}
\definecolor{dblue}{RGB}{150,50,150}
%\usepackage[colorlinks=true, linkcolor=dblue, urlcolor=cyan]{hyperref}
\usepackage{hyperref}
%\usepackage{cleveref}

%\usepackage{crossreftools}

%\addto\captionsngerman{%
 \renewcommand{\abstractname}{Vorbemerkung}
%}

\title{Symmetrien und Erhaltungs-Sätze}
\begin{document}
\author{Hugo Meder}
\begin{titlepage}
\maketitle
\end{titlepage}
\tableofcontents

\begin{abstract}
In diesem Teil geht es darum, Symmetrien physikalischer Systeme zu untersuchen. Das Noethersche Theorem schafft einen Zusammenhang
zwischen bestimmten Symmetrien und zeitlich unveränderlichen Größen, wie Energie, Impuls usw.
\end{abstract}

\begin{section}{Symmetrien}
\begin{subsection}{Transformationen}
Wir betrachten für ein bestimmtes physikalisches System seinen Konfigurations-Raum $K$. Im Falle von Systemen mit verallgemeinerten Koordinate sind dies $alle$ denkbaren Koordinaten $q=(q_1,q_2,...,q_n)$ wobei $n$ die Anzahl der Koordinaten ist\footnote{Hier schließen wir auch Systeme von Feldern, z.B elektrische und Magnetische Felder nicht aus. In solchen Fällen ist der Konfigurations-Raum
dann allerdings unendlich-dimensional.}. Weiter sei $\tilde{E}$ der Ereignis-Raum. Dieser besteht aus Paaren
\begin{eqnarray*}
(q,t) &\in& \tilde{E}
\end{eqnarray*}
Ereignisse sind mathematische Objekte (Elemente der Menge $\tilde{E}$), nicht Aussagen. Wenn wir aber sagen ein Ereignis
$(p,t)$ habe stattgefunden oder wird stattfinden, dann bedeutet dies, dass sich das System zum Zeitpunkt $t$ eine Konfiguration $q$
hat oder hatte. Dies ist eine Aussage und wir schreiben
\begin{eqnarray*}
q(t) &=& q
\end{eqnarray*}
Eine Transformation $\tilde{T}$ ist eine Abbildung, die Ereignissen $e\in\tilde{E}$ andere Ereignisse $\tilde{T}(e)\in\tilde{E}$
zuordnet. 
\begin{eqnarray*}
\tilde{T}: \tilde{E} &\rightarrow&\tilde{E}\\
(q,t) &\mapsto& ( q^\prime,t^\prime)=\left(\tilde{T}_q(q,t), \tilde{T}_t(q,t)\right)
\end{eqnarray*}
Man beachte hier, dass eine Transformation nicht nur die Konfiguration $p$ sondern auch die dem Ereignis zugeordnete Zeit $t$ verändern kann.\footnote{In der Praxis sind nur Transformationen zu betrachten, deren beide Teile $T_q$ und $T_t$ glatt sind, dass heißt, stetig ableitbar sind.}
Um ein anschaulicheres Bild zu erhalten, betrachten wir Systeme, deren Konfigurationsraum aus einer Anzahl $n$ von Raumkoordinaten $x_k$ besteht.
\begin{eqnarray*}
\tilde{T}: \tilde{E} &\rightarrow&\tilde{E}\\
(x_k,t) &\mapsto& ( x_k^\prime,t^\prime)=\left(T_x(x_k,t), T_t(x_k,t)\right)
\end{eqnarray*}
In diesem Falle besteht eine Transformation in einer Veränderung der Koordinaten $x_k$,und der Zeit $t$.
Wichtige Beispiele sind:
\begin{itemize}
	\item räumliche Verschiebungen ($x_k^\prime=x_k+y, t^\prime=t$)
	\item zeitliche Verschiebungen ($x_k^\prime=x_k, t^\prime=t+\Delta t$)
	\item Raumspiegelung ($x_k^\prime=-x_k, t^\prime=t$)
	\item räumlich Drehungen ($x_k^\prime=Rx_k, t^\prime=t$)
	\item Zeitumkehr ($x_k^\prime=x_k, t^\prime=-t$)
\end{itemize}
Spiegelungen an Ebenen erhält man indem zuerst eine Raumspiegelung durchgeführt wird und dann um eine beliebige Achse um 180° gedreht wird. Dies ist dann das gleiche wie eine Spiegelung an der zur Achse senkrechten Ebene.

Transformationen können nacheinander auf Ereignisse angewandt werden. Es seien $\tilde{T}^1$ und $\tilde{T}^2$ Transformationen, bezeichnen
wir mit $\tilde{T}^1\circ \tilde{T}^2$ die Transformation
\begin{eqnarray*}
\tilde{T}^1\circ \tilde{T}^2: \tilde{E} &\rightarrow&\tilde{E}\\
(q,t) &\mapsto& ( q^\prime,t^\prime)=\tilde{T}^1\left(\tilde{T}^2(q_k,t)\right)=\left(\tilde{T}^1_q\left(\tilde{T}^2(q,t)\right), \tilde{T}^1_t\left(\tilde{T}^2(q,t)\right)\right)
\end{eqnarray*}
Eine besonders einfache Transformation $\tilde{I}$ lässt alles unverändert und wird Identität genannt:
\begin{eqnarray*}
\tilde{I}: \tilde{E} &\rightarrow&\tilde{E}\\
(q,t) &\mapsto& ( q^\prime,t^\prime)=( q,t)
\end{eqnarray*}
Gibt es zu einer Transformation $\tilde{T}$ eine Transformation $\tilde{T}^\prime$ mit der Eigenschaft dass $\tilde{T}^\prime\circ\tilde{T}=\tilde{I}$ dann nennen wir $\tilde{T}^\prime$ die zu $\tilde{T}$ inverse Transformation und gezeichnen sie mit $\tilde{T}^{-1}$.

Eine Menge $G$ von Transformationen heißt \textbf{Transformations-Gruppe}, wenn folgendes gilt:
\begin{itemize}
	\item $\tilde{I}$ ist Element von $G$.
	\item ist $\tilde{T}^1$ und $\tilde{T}^2$ in $G$ enthalten, dann ist auch $\tilde{T}^1\circ\tilde{T}^2$ in $G$ enthalten.
	\item ist $\tilde{T}$ in $G$ enthalten, dann ist auch die Inverse $\tilde{T}^{-1}$ in $G$ enthalten.
\end{itemize}
Jede der oben aufgeführten Transformationen bilden je eine Transformations-Gruppe. Die beiden diskreten Transformation, die  Raumspiegelung und der Zeitumkehr bestehen jeweils aus zwei Elementen, der Transformation und der Identität $\tilde{I}$
Verknüpft man diese mittels $\circ$ miteinander, dann erhält man eine einzige grössere Transformations-Gruppe.
\end{subsection}

\begin{subsection}{Physikalische Symmetrien}
Eine dynamische Zeitentwicklung eines Systems ordnet jeder Zeit $t$ eine Konfiguration $q_k(t)$ zu. Diese Zeitentwicklung soll glatt sein, so
dass aus der Zeitentwicklung auch die Zeitableitungen $\dot{q}_k(t)$ und $\ddot{q}_k(t)$ hervorgehen. Erfüllt eine Zeitentwicklung die Bewegungsgleichungen, dann sprechen wir von einer physikalischen Zeitentwicklung. Zeitentwicklungen können mittels einer Transformation $\tilde{T}$ in andere Zeitentwicklungen überführt werden. Eine Zeitentwicklung kann als eine Menge von Ereignissen $(q(t),t)$ aufgefasst werden. Diese können Transformiert werden und gehen dadurch in $(T_q(q_k(t)),T_t(t))$ über. Daraus lässt sich wieder eine Zeitentwicklung $q^\prime(t)$ gewinnen. Um diesen Vorgang zu erläutern, betrachten wir die im letzten Abschnitt beschriebenen Transformationen:
\begin{itemize}
	\item räumliche Verschiebungen: $x_k^\prime(t) = x_k(t)+y$
	\item zeitliche Verschiebungen $x_k^\prime(t) = x_k(t-\Delta t)$
	\item Raumspiegelung $x_k^\prime(t) = -x_k(t)$
	\item räumlich Drehungen $x_k^\prime(t) = Rx_k(t)$
	\item Zeitumkehr $x_k^\prime(t) = x_k(-t)$
\end{itemize}
Eine Transformation $\tilde{T}$ nennen wir eine \textbf{physikalische Symmetrie} falls sie \textbf{alle} physikalische 
Zeitentwicklungen in physikalische Zeitentwicklungen abbildet. Bildet eine Menge physikalischer Symmetrien eine Transformations-Gruppe,
dann nennen wir diese eine \textbf{physikalische Symmetrie-Gruppe}. Für viele wichtige Systeme sind die im letzten Abschnitt aufgeführten Transformations-Gruppen auch physikalische Symmetrie-Gruppen.
\end{subsection}

\begin{subsection}{Kontinuierliche Transformationen}
Viele Transformations-Gruppen sind Lie-Gruppen. Diese zeichnen sich dadurch aus, dass es Umgebungen der Identität $\tilde{I}$ gibt, die deren Elemente durch reelle Zahlen parametrisiert werden können. Sind dies bezüglich einem System physikalische Symmetrie-Gruppen, dann spricht man auch von kontinuierlichen physikalischen Symmetrien. Greift man einen der Parameter heraus, dann erhält man eine einparametrige
Transformations-Gruppe die auch eine einparametrige
Symmetrie-Gruppe sein kann. Für einparametrige Transformations-Gruppen bezeichnen wir die Transformationen mit $\tilde{T}(\epsilon)$ wobei $\epsilon$ den Parameter bezeichnet: Es gelte $\tilde{T}(0)=\tilde{I}$ und somit $\tilde{T}(0)(q_k,t)=(q_k,t)$.
\begin{eqnarray*}
\tilde{T}(\epsilon): \tilde{E} &\rightarrow&\tilde{E}\\
(q_k,t) &\mapsto& ( q_k^\prime,t^\prime)=\left(T_q(\epsilon,q_k,t), T_t(\epsilon,q_k.t)\right)
\end{eqnarray*}
Ist $\epsilon$ sehr klein, dann haben wir
\begin{eqnarray*}
T_q(\epsilon,q_k,t) &=& q_k+\vartheta_q(q_k,t)\epsilon\\
T_t(\epsilon,q_k.t) &=& t+\vartheta_t(q_k,t)\epsilon
\end{eqnarray*}
wobei
\begin{eqnarray*}
\vartheta_q(q_k,t) &=& \left(\frac{\partial T_q(\epsilon,q_k,t)}{\partial \epsilon}\right|_{\epsilon=0}\\
\vartheta_t(q_k,t) &=& \left(\frac{\partial T_t(\epsilon,q_k,t)}{\partial \epsilon}\right|_{\epsilon=0}
\end{eqnarray*}
Dies kann auch so geschrieben werden:
\begin{eqnarray*}
\tilde{T}(\epsilon)(q_k,t) &=& (q_k,t)+(\vartheta_q(q_k,t),\vartheta_t(q_k,t))\epsilon
\end{eqnarray*}
Für eine infinitesimale Transformation $\tilde{T}(\epsilon)$ mit infinitesimalem $\epsilon$ lässt sich die Transformation
einer Zeitentwicklungen (bis zur ersten Ordnung in $\epsilon$) direkt hinschreiben:
\begin{eqnarray*}
q_k^\prime(t+\vartheta_t(q_k,t)\epsilon) &=& q_k(t)+\vartheta_{q_k}(q_k,t)\epsilon\\
q_k^\prime(t) &=& q_k(t-\vartheta_t(q_k,t)\epsilon)+\vartheta_{q_k}(q_k,t-\vartheta_t(q_k,t)\epsilon)\epsilon\\
&=& q_k(t-\vartheta_t(q_k,t)\epsilon)+\vartheta_{q_k}(q_k,t)\epsilon\\
&=& q_k(t)-\dot{q}_k(t)\vartheta_t(q_k,t)\epsilon+\vartheta_{q_k}(q_k,t)\epsilon\\
&=& q_k(t)+\left(-\dot{q}_k(t)\vartheta_t(q_k(t),t)+\vartheta_{q_k}(q_k(t),t)\right)\epsilon
\end{eqnarray*}
Ebenso lässt sich $\dot{q}^\prime(t)$ hinschreiben:
\begin{eqnarray*}
\dot{q}_k^\prime(t) &=& \frac{d}{dt}\left(q_k(t)-\dot{q}_k(t)\vartheta_t(q_k(t),t)\epsilon+\vartheta_{q_k}(q_k(t),t)\epsilon\right)\\
&=& \dot{q}_k(t)\\
&&-\left(\ddot{q}_k(t)\vartheta_t(q_k(t),t)+\dot{q}_k(t)\sum_l\frac{\partial\vartheta_t(q_k(t),t)}{\partial q_l}\dot{q}_l(t)+\dot{q}_k(t)\frac{\partial\vartheta_t(q_k(t),t)}{\partial t}\right)\epsilon\\
&&+\left(\sum_l\frac{\partial\vartheta_{q_k}(q_k(t),t)}{\partial q_l}\dot{q}_l+\frac{\partial\vartheta_{q_k}(q_k(t),t)}{\partial t}\right)\epsilon\\
&=& \dot{q}_k(t)\left(1-\frac{d\vartheta_t(q_k(t),t)}{dt}\epsilon\right)-\ddot{q}_k(t)\vartheta_t(q_k(t),t)\epsilon
+\frac{d\vartheta_{q_k}(q_k(t),t)}{dt}\epsilon
\end{eqnarray*}
Wir definieren:
\begin{eqnarray}
\label{by_d_epsilon_def}
\frac{dq_k}{d\epsilon}(\dot{q},q) &=& \left(\frac{\partial}{\partial\epsilon}q^\prime_k(t)\right|_{\epsilon=0}
=-\dot{q}_k(t)\vartheta_t(q_k(t),t)+\vartheta_{q_k}(q_k(t),t)\cr
\frac{d\dot{q}_k}{d\epsilon}(\ddot{q},\dot{q},q) &=& \left(\frac{\partial}{\partial\epsilon}\dot{q}^\prime_k(t)\right|_{\epsilon=0}
=-\dot{q}_k(t)\frac{d\vartheta_t(q_k(t),t)}{dt}-\ddot{q}_k(t)\vartheta_t(q_k(t),t)
+\frac{d\vartheta_{q_k}(q_k(t),t)}{dt}
\end{eqnarray}
\end{subsection}

\begin{subsection}{Infinitesimale Transformationen des Phasen-Raums}
\begin{subsubsection}{Problemstellung}
Ziel ist es hier infinitesimale Transformationen in die Form
\begin{eqnarray*}
q^\prime_k &=& q_k+G_k(p,q)\epsilon\\
p^\prime_k &=& p_k+F_k(p,q)\epsilon
\end{eqnarray*}
zu bringen.
Zur Illustration dient eine infinitesimale Rotation um die $z$-Achse. Die konkrete Form der Lagrange-Funtion $L(\dot{q},q)$ lassen wir zunächst offen. Es gilt:
\begin{eqnarray*}
q_k^\prime &=& q_k+\sum_m\sum_l\varepsilon_{kml}n_mq_l\epsilon
\end{eqnarray*}
Ist $n$ eine Einheitsvektor in $3$-Richtung bedeutet dies:
\begin{eqnarray*}
q_1^\prime &=& q_1+\sum_l\varepsilon_{13l}q_l\epsilon\\
&=& q_1-q_2\epsilon\\
q_2^\prime &=& q_2+\sum_l\varepsilon_{23l}q_l\epsilon\\
&=& q_2+q_1\epsilon\\
q_3^\prime &=& q_3+\sum_l\varepsilon_{33l}q_l\epsilon\\
&=& q_3
\end{eqnarray*}
Wir erhalten weiter
\begin{eqnarray*}
\frac{dq_k}{d\epsilon} &=& \sum_m\sum_l\varepsilon_{kml}n_mq_l\\
\frac{d\dot{q}_k}{d\epsilon} &=& \sum_m\sum_l\varepsilon_{kml}n_m\dot{q}_l
\end{eqnarray*}

\end{subsubsection}

\begin{subsubsection}{Infinitesimale Transformationen der konjugierten Impulse}
\begin{eqnarray*}
p_k^\prime(t) &=& \frac{\partial L}{\partial \dot{q}_k}(\dot{q}^\prime_k(t),q^\prime_k(t))
=\frac{\partial L}{\partial \dot{q}_k}\left(\dot{q}_k(t)+\frac{d\dot{q}_k}{d\epsilon}\epsilon,q_k(t)+\frac{dq_k}{d\epsilon}\epsilon\right)\\
&=& \frac{\partial L}{\partial \dot{q}_k}\left(\dot{q}_k(t),q_k(t)\right)
+\sum_l\frac{\partial^2 L}{\partial \dot{q}_k\partial \dot{q}_l}\left(\dot{q}_k(t),q_k(t)\right)\frac{d\dot{q}_l}{d\epsilon}\epsilon
+\sum_l\frac{\partial^2 L}{\partial \dot{q}_k\partial q_l}\left(\dot{q}_k(t),q_l(t)\right)\frac{dq_l}{d\epsilon}\epsilon\\
&=& p_k(\dot{q},q)
+\left(\sum_l\frac{d\dot{q}_l}{d\epsilon}\frac{\partial}{\partial \dot{q}_k}\frac{\partial L}{\partial \dot{q}_l}\left(\dot{q}_k(t),q_k(t)\right)
+\sum_l\frac{dq_l}{d\epsilon}\frac{\partial}{\partial \dot{q}_k}\frac{\partial L}{\partial q_l}\left(\dot{q}_k(t),q_l(t)\right)\right)\epsilon\\
&=& p_k(\dot{q},q)
+\left(\sum_l\frac{d\dot{q}_l}{d\epsilon}\frac{\partial}{\partial \dot{q}_k}p_l\left(\dot{q}_k(t),q_k(t)\right)
+\sum_l\frac{dq_l}{d\epsilon}\frac{\partial}{\partial \dot{q}_k}\frac{\partial L}{\partial q_l}\left(\dot{q}_k(t),q_l(t)\right)\right)\epsilon\\
\frac{dp_k}{d\epsilon}\left(\dot{q}_k(t),q_l(t)\right) &=&\sum_l\frac{d\dot{q}_l}{d\epsilon}\frac{\partial}{\partial \dot{q}_k}p_l\left(\dot{q}_k(t),q_k(t)\right)
+\sum_l\frac{dq_l}{d\epsilon}\frac{\partial}{\partial \dot{q}_k}\frac{\partial L}{\partial q_l}\left(\dot{q}_k(t),q_l(t)\right)
\end{eqnarray*}
\begin{paragraph}{Beispiel Rotation}hh
\begin{eqnarray*}
p_k^\prime(t) &=& p_k(\dot{q},q)
+\left(\sum_l\frac{d\dot{q}_l}{d\epsilon}\frac{\partial}{\partial \dot{q}_k}p_l\left(\dot{q}_k(t),q_k(t)\right)
+\sum_l\frac{dq_l}{d\epsilon}\frac{\partial}{\partial \dot{q}_k}\frac{\partial L}{\partial q_l}\left(\dot{q}_k(t),q_l(t)\right)\right)\epsilon\\
&=& p_k(\dot{q},q)+\sum_l\sum_{rs}\varepsilon_{lrs}n_r\dot{q}_s\frac{\partial}{\partial \dot{q}_k}p_l\left(\dot{q}_k(t),q_k(t)\right)\epsilon\\
&&+\sum_l\sum_{rs}\varepsilon_{lrs}n_rq_s\frac{\partial}{\partial \dot{q}_k}\frac{\partial L}{\partial q_l}\left(\dot{q}_k(t),q_l(t)\right)\epsilon\\
\frac{dp_k}{d\epsilon}\left(\dot{q}_k(t),q_l(t)\right) &=&
\sum_{rs}\varepsilon_{lrs}n_r\left(\dot{q}_s\frac{\partial}{\partial \dot{q}_k}p_l\left(\dot{q}_k(t),q_k(t)\right)
+q_s\frac{\partial}{\partial \dot{q}_k}\frac{\partial L}{\partial q_l}\left(\dot{q}_k(t),q_l(t)\right)\right)
\end{eqnarray*}
\end{paragraph}
\end{subsubsection}

\begin{subsubsection}{Annahme: \titlemath{$\vartheta_t=0$} und \titlemath{$\frac{\partial\vartheta_q}{\partial t}=0$}}
Dann wird aus \ref{by_d_epsilon_def}:
\begin{eqnarray*}
\frac{dq_k}{d\epsilon}(q) &=& \vartheta_{qk}(q)\cr
\frac{d\dot{q}_k}{d\epsilon}(\dot{q},q) &=&\frac{d\vartheta_{qk}(q_k(t),t)}{dt}
=\sum_l\frac{\partial\vartheta_{qk}}{\partial q_l}(q)\dot{q}_l
\end{eqnarray*}
Dies hat zur Folge, dass
\begin{eqnarray*}
\frac{\partial}{\partial \dot{q}_k}\frac{d\dot{q}_l}{d\epsilon} &=& \frac{\partial}{\partial \dot{q}_k}\sum_r\frac{\partial\vartheta_{ql}}{\partial q_r}(q)\dot{q}_r=\frac{\partial\vartheta_{ql}}{\partial q_k}(q)\\
\frac{\partial}{\partial \dot{q}_k}\frac{dq_l}{d\epsilon} &=& 0
\end{eqnarray*}
und somit
\begin{eqnarray*}
\frac{dp_k}{d\epsilon}\left(\dot{q}_k(t),q_l(t)\right) &=&\sum_l\frac{d\dot{q}_l}{d\epsilon}\frac{\partial}{\partial \dot{q}_k}p_l\left(\dot{q}_k(t),q_k(t)\right)
+\sum_l\frac{dq_l}{d\epsilon}\frac{\partial}{\partial \dot{q}_k}\frac{\partial L}{\partial q_l}\left(\dot{q}_k(t),q_l(t)\right)\\
&=&\frac{\partial}{\partial \dot{q}_k}\sum_l\left(\frac{d\dot{q}_l}{d\epsilon}p_l\left(\dot{q}_k(t),q_k(t)\right)
+\frac{dq_l}{d\epsilon}\frac{\partial L}{\partial q_l}\left(\dot{q}_k(t),q_l(t)\right)\right)\\
&&-\sum_l\frac{\partial\vartheta_{ql}}{\partial q_k}(q)p_l\left(\dot{q}_k(t),q_k(t)\right)
\end{eqnarray*}
\begin{paragraph}{Beispiel Rotation}
In diesem Falle gilt:
\begin{eqnarray*}
\vartheta_{qk}(q) &=& \sum_{rs}\varepsilon_{krs}n_rq_s
\end{eqnarray*}
Und somit:
\begin{eqnarray*}
\frac{dp_k}{d\epsilon}\left(\dot{q}_k(t),q_l(t)\right) &=&\frac{\partial}{\partial \dot{q}_k}\sum_l\left(\frac{d\dot{q}_l}{d\epsilon}p_l\left(\dot{q}_k(t),q_k(t)\right)
+\frac{dq_l}{d\epsilon}\frac{\partial L}{\partial q_l}\left(\dot{q}_k(t),q_l(t)\right)\right)\\
&&-\sum_l\frac{\partial\vartheta_{ql}}{\partial q_k}(q)p_l\left(\dot{q}_k(t),q_k(t)\right)\\
&=&\frac{\partial}{\partial \dot{q}_k}\sum_l\left(\sum_{rs}\varepsilon_{lrs}n_r\dot{q}_sp_l\left(\dot{q}_k(t),q_k(t)\right)
+\sum_{rs}\varepsilon_{lrs}n_rq_s\frac{\partial L}{\partial q_l}\left(\dot{q}_k(t),q_l(t)\right)\right)\\
&&-\sum_l\sum_r\varepsilon_{lrk}n_rp_l\left(\dot{q}_k(t),q_k(t)\right)
\end{eqnarray*}
\end{paragraph}
\end{subsubsection}
\end{subsection}

\begin{subsection}{Scratch}
Wir untersuchen nun die Lagrange-Bewegungsgleichungen für $q^\prime_k(t)=q_k(t)+\frac{dq_k(t)}{d\epsilon}\epsilon$. Für den kanonisch konjugierten Impuls $p^\prime(t)$ erhalten wir:
\begin{eqnarray*}
p_k^\prime(t) &=& \frac{\partial L}{\partial \dot{q}_k}(\dot{q}^\prime_k(t),q^\prime_k(t))
=\frac{\partial L}{\partial \dot{q}_k}\left(\dot{q}_k(t)+\frac{d\dot{q}_k}{d\epsilon}\epsilon,q_k(t)+\frac{dq_k}{d\epsilon}\epsilon\right)\\
&=& \frac{\partial L}{\partial \dot{q}_k}\left(\dot{q}_k(t),q_k(t)\right)
+\sum_l\frac{\partial^2 L}{\partial \dot{q}_k\partial \dot{q}_l}\left(\dot{q}_k(t),q_k(t)\right)\frac{d\dot{q}_l}{d\epsilon}\epsilon
+\sum_l\frac{\partial^2 L}{\partial \dot{q}_k\partial q_l}\left(\dot{q}_k(t),q_l(t)\right)\frac{dq_l}{d\epsilon}\epsilon
\end{eqnarray*}
Und ebenso für
\begin{eqnarray*}
\frac{\partial L}{\partial q_k}(\dot{q}^\prime_k(t),q^\prime_k(t))
&=&\frac{\partial L}{\partial q_k}\left(\dot{q}_k(t)+\frac{d\dot{q}_k}{d\epsilon}\epsilon,q_k(t)+\frac{dq_k}{d\epsilon}\epsilon\right)\\
&=& \frac{\partial L}{\partial q_k}\left(\dot{q}_k(t),q_k(t)\right)\\
&&+\sum_l\frac{\partial^2 L}{\partial q_k\partial \dot{q}_l}\left(\dot{q}_k(t),q_k(t)\right)\frac{d\dot{q}_l}{d\epsilon}\epsilon
+\sum_l\frac{\partial^2 L}{\partial q_k\partial q_l}\left(\dot{q}_k(t),q_l(t)\right)\frac{dq_l}{d\epsilon}\epsilon
\end{eqnarray*}
Mit der Definition
\begin{eqnarray*}
\frac{dL}{d\epsilon}&=&\sum_l\frac{\partial L}{\partial \dot{q}_l}\left(\dot{q}_k(t),q_k(t)\right)\frac{d\dot{q}_l}{d\epsilon}
+\sum_l\frac{\partial L}{\partial q_l}\left(\dot{q}_k(t),q_l(t)\right)\frac{dq_l}{d\epsilon}
\end{eqnarray*}
verkürzen sich beide Ausdrücke:
\begin{eqnarray*}
p_k^\prime(t) &=& \frac{\partial L}{\partial \dot{q}_k}\left(\dot{q}_k(t),q_k(t)\right)
+\frac{\partial}{\partial \dot{q}_k}\frac{dL}{d\epsilon}\epsilon\\
\frac{\partial L}{\partial q_k}(\dot{q}^\prime_k(t),q^\prime_k(t))
&=& \frac{\partial L}{\partial q_k}\left(\dot{q}_k(t),q_k(t)\right)
+\frac{\partial}{\partial q_k}\frac{dL}{d\epsilon}\epsilon
\end{eqnarray*}
Wir nehmen nun an, dass $q_k(t)$ die Bewegungsgleichungen
\begin{eqnarray*}
\frac{d}{dt}\frac{\partial L}{\partial \dot{q}_k}\left(\dot{q}_k(t),q_k(t)\right)-\frac{\partial L}{\partial q_k}\left(\dot{q}_k(t),q_k(t)\right) &=& 0
\end{eqnarray*}
erfüllt. Damit auch $q^\prime_k(t)$
\begin{eqnarray*}
\frac{d}{dt}p^\prime_k(t)-\frac{\partial L}{\partial q_k}(\dot{q}^\prime_k(t),q^\prime_k(t)) &=& 0
\end{eqnarray*}
muss
\begin{eqnarray*}
\frac{d}{dt}\frac{\partial}{\partial \dot{q}_k}\frac{dL}{d\epsilon}\epsilon-\frac{\partial}{\partial q_k}\frac{dL}{d\epsilon}\epsilon &=&0\\
\frac{d}{dt}\frac{\partial}{\partial \dot{q}_k}\frac{dL}{d\epsilon}-\frac{\partial}{\partial q_k}\frac{dL}{d\epsilon} &=&0
\end{eqnarray*}
gelten. Unter der Bedingung, dass $q_k(t)$ die Bewegungsgleichungen erfüllt gilt:
\begin{eqnarray*}
\frac{dL}{d\epsilon}&=&\sum_l\frac{\partial L}{\partial \dot{q}_l}\left(\dot{q}_k(t),q_k(t)\right)\left(-\dot{q}_l(t)\frac{d\vartheta_t(q_k(t),t)}{dt}
+\frac{d\vartheta_{q_l}(q_k(t),t)}{dt}\right)\\
&&+\sum_l\left(\frac{d}{dt}\frac{\partial L}{\partial \dot{q}_l}\left(\dot{q}_k(t),q_l(t)\right)\right)\vartheta_{q_l}(q_k(t),t)\\
&&+\frac{dL}{dt}\left(\dot{q}_k(t),q_k(t)\right)\vartheta_t(q_k(t),t)\\
&=&-\left(\sum_l\frac{\partial L}{\partial \dot{q}_l}\left(\dot{q}_k(t),q_k(t)\right)\dot{q}_l(t)\right)\frac{d\vartheta_t(q_k(t),t)}{dt}\\
&&+\frac{d}{dt}\left(\sum_l\frac{\partial L}{\partial \dot{q}_l}\left(\dot{q}_k(t),q_l(t)\right)\vartheta_{q_l}(q_k(t),t)\right)\\
&&+\frac{dL}{dt}\left(\dot{q}_k(t),q_k(t)\right)\vartheta_t(q_k(t),t)
\end{eqnarray*}
.............................................................
\begin{eqnarray*}
\frac{dL}{d\epsilon}&=&\sum_l\frac{\partial L}{\partial \dot{q}_l}(\dot{q}_k,q_k)\frac{d\dot{q}_l}{d\epsilon}(\dot{q}_k,q_k)
+\sum_l\frac{\partial L}{\partial q_l}(\dot{q}_k,q_k)\frac{dq_l}{d\epsilon}(\ddot{q}_k,\dot{q}_k,q_k)\\
\frac{\partial}{\partial q_k}\frac{dL}{d\epsilon}&=&\sum_l\frac{\partial^2 L}{\partial \dot{q}_l\partial q_k}(\dot{q}_k,q_k)\frac{d\dot{q}_l}{d\epsilon}(\ddot{q}_k,\dot{q}_k,q_k)
+\sum_l\frac{\partial^2 L}{\partial q_l\partial q_k}(\dot{q}_k,q_k)\frac{dq_l}{d\epsilon}(\dot{q}_k,q_k)\\
&&+\sum_l\frac{\partial L}{\partial \dot{q}_l}(\dot{q}_k,q_k)\frac{\partial}{\partial q_k}\frac{d\dot{q}_l}{d\epsilon}(\ddot{q}_k,\dot{q}_k,q_k)
+\sum_l\frac{\partial L}{\partial q_l}(\dot{q}_k,q_k)\frac{\partial}{\partial q_k}\frac{dq_l}{d\epsilon}(\dot{q}_k,q_k)\\
&=&\frac{d}{d\epsilon}\frac{\partial L}{\partial q_k}(\dot{q}_k,q_k)\\
&&+\sum_l\frac{\partial L}{\partial \dot{q}_l}(\dot{q}_k,q_k)\frac{\partial}{\partial q_k}\frac{d\dot{q}_l}{d\epsilon}(\ddot{q}_k,\dot{q}_k,q_k)
+\sum_l\frac{\partial L}{\partial q_l}(\dot{q}_k,q_k)\frac{\partial}{\partial q_k}\frac{dq_l}{d\epsilon}(\dot{q}_k,q_k)\\
&=&\frac{d}{d\epsilon}\frac{\partial L}{\partial q_k}(\dot{q}_k,q_k)\\
&&+\sum_l\frac{\partial L}{\partial \dot{q}_l}(\dot{q}_k,q_k)\frac{\partial}{\partial q_k}\frac{d\dot{q}_l}{d\epsilon}(\ddot{q}_k,\dot{q}_k,q_k)
+\sum_l\frac{d}{dt}\frac{\partial L}{\partial \dot{q}_l}(\dot{q}_k,q_k)\frac{\partial}{\partial q_k}\frac{dq_l}{d\epsilon}(\dot{q}_k,q_k)\\
&=&\frac{d}{d\epsilon}\frac{\partial L}{\partial q_k}(\dot{q}_k,q_k)+\frac{d}{dt}\sum_l\frac{\partial L}{\partial \dot{q}_l}(\dot{q}_k,q_k)\frac{\partial}{\partial q_k}\frac{dq_l}{d\epsilon}(\dot{q}_k,q_k)\\
\frac{\partial}{\partial \dot{q}_k}\frac{dL}{d\epsilon}&=&\frac{d}{d\epsilon}\frac{\partial L}{\partial \dot{q}_k}(\dot{q}_k,q_k)\\
&&+\sum_l\frac{\partial L}{\partial \dot{q}_l}(\dot{q}_k,q_k)\frac{\partial}{\partial \dot{q}_k}\frac{d\dot{q}_l}{d\epsilon}(\ddot{q}_k,\dot{q}_k,q_k)
+\sum_l\frac{\partial L}{\partial q_l}(\dot{q}_k,q_k)\frac{\partial}{\partial \dot{q}_k}\frac{dq_l}{d\epsilon}(\dot{q}_k,q_k)\\
&=&\frac{d}{d\epsilon}\frac{\partial L}{\partial \dot{q}_k}(\dot{q}_k,q_k)\\
&&+\sum_l\frac{\partial L}{\partial \dot{q}_l}(\dot{q}_k,q_k)\frac{\partial}{\partial \dot{q}_k}\frac{d\dot{q}_l}{d\epsilon}(\ddot{q}_k,\dot{q}_k,q_k)
+\sum_l\frac{d}{dt}\frac{\partial L}{\partial \dot{q}_l}(\dot{q}_k,q_k)\frac{\partial}{\partial \dot{q}_k}\frac{dq_l}{d\epsilon}(\dot{q}_k,q_k)\\
&=&\frac{d}{d\epsilon}\frac{\partial L}{\partial \dot{q}_k}(\dot{q}_k,q_k)
+\frac{d}{dt}\sum_l\frac{\partial L}{\partial \dot{q}_l}(\dot{q}_k,q_k)\frac{\partial}{\partial \dot{q}_k}\frac{dq_l}{d\epsilon}(\dot{q}_k,q_k)\\
\frac{d}{dt}\frac{\partial}{\partial \dot{q}_k}\frac{dL}{d\epsilon}
&=&\frac{d}{dt}\frac{d}{d\epsilon}\frac{\partial L}{\partial \dot{q}_k}(\dot{q}_k,q_k)\\
&&+\frac{d^2}{dt^2}\sum_l\frac{\partial L}{\partial \dot{q}_l}(\dot{q}_k,q_k)\frac{\partial}{\partial \dot{q}_k}\frac{dq_l}{d\epsilon}(\dot{q}_k,q_k)\\
&=&\frac{d}{dt}\left(\sum_l\frac{\partial^2 L}{\partial \dot{q}_k\partial \dot{q}_l}(\dot{q}_k,q_k)\frac{d\dot{q}_l}{d\epsilon}
+\sum_l\frac{\partial^2 L}{\partial \dot{q}_k\partial q_l}(\dot{q}_k,q_k)\frac{dq_l}{d\epsilon}\right)\\
&&+\frac{d^2}{dt^2}\sum_l\frac{\partial L}{\partial \dot{q}_l}(\dot{q}_k,q_k)\frac{\partial}{\partial \dot{q}_k}\frac{dq_l}{d\epsilon}(\dot{q}_k,q_k)
\end{eqnarray*}
..........................................................................
\begin{eqnarray*}
\dot{p}^\prime_k &=& -\frac{\partial}{\partial q_k}H(p^\prime,q^\prime)\\
\dot{q}^\prime_k &=& \frac{\partial}{\partial p_k}H(p^\prime,q^\prime)\\
\dot{p}_k+\frac{d\dot{p}_k}{d\epsilon}\epsilon &=& -\frac{\partial}{\partial q_k}H\left(p+\epsilon\frac{dp}{d\epsilon},q+\epsilon\frac{dq}{d\epsilon}\right)\\
\dot{q}_k+\frac{d\dot{q}_k}{d\epsilon}\epsilon &=& \frac{\partial}{\partial q_k}H\left(p+\epsilon\frac{dp}{d\epsilon},q+\epsilon\frac{dq}{d\epsilon}\right)\\
\dot{p}_k+\frac{d\dot{p}_k}{d\epsilon}\epsilon &=& -\frac{\partial}{\partial q_k}H\left(p,q\right)
-\epsilon\sum_l\frac{\partial^2}{\partial q_k\partial p_l}H\left(p,q\right)\frac{dp_l}{d\epsilon}
-\epsilon\sum_l\frac{\partial^2}{\partial q_k\partial q_l}H\left(p,q\right)\frac{dq_l}{d\epsilon}\\
\dot{q}_k+\frac{d\dot{q}_k}{d\epsilon}\epsilon &=& \frac{\partial}{\partial p_k}H\left(p,q\right)
+\epsilon\sum_l\frac{\partial^2}{\partial p_k\partial p_l}H\left(p,q\right)\frac{dp_l}{d\epsilon}
+\epsilon\sum_l\frac{\partial^2}{\partial p_k\partial q_l}H\left(p,q\right)\frac{dq_l}{d\epsilon}\\
\frac{d\dot{p}_k}{d\epsilon} &=& 
-\sum_l\frac{\partial^2}{\partial q_k\partial p_l}H\left(p,q\right)\frac{dp_l}{d\epsilon}
-\sum_l\frac{\partial^2}{\partial q_k\partial q_l}H\left(p,q\right)\frac{dq_l}{d\epsilon}\\
\frac{d\dot{q}_k}{d\epsilon} &=& 
+\sum_l\frac{\partial^2}{\partial p_k\partial p_l}H\left(p,q\right)\frac{dp_l}{d\epsilon}
+\sum_l\frac{\partial^2}{\partial p_k\partial q_l}H\left(p,q\right)\frac{dq_l}{d\epsilon}\\
\frac{d\dot{p}_k}{d\epsilon} &=& 
\sum_l\frac{\partial}{\partial p_l}\dot{p}_k\frac{dp_l}{d\epsilon}
+\sum_l\frac{\partial}{\partial q_l}\dot{q}_k\frac{dq_l}{d\epsilon}\\
\frac{d\dot{q}_k}{d\epsilon} &=& 
+\sum_l\frac{\partial^2}{\partial p_k\partial p_l}H\left(p,q\right)\frac{dp_l}{d\epsilon}
+\sum_l\frac{\partial^2}{\partial p_k\partial q_l}H\left(p,q\right)\frac{dq_l}{d\epsilon}\\
\end{eqnarray*}
Poisson-Klammer:
\begin{eqnarray*}
\{A(p,q),B(p,q)\} &=& \sum_k\left(\frac{\partial A}{\partial q_k}(p,q)\frac{\partial B}{\partial p_k}(p,q)
-\frac{\partial A}{\partial p_k}(p,q)\frac{\partial B}{\partial q_k}(p,q)\right)\\
\{A(p,q),H(p,q)\} &=& \sum_k\left(\frac{\partial A}{\partial q_k}(p,q)\frac{\partial H}{\partial p_k}(p,q)
-\frac{\partial A}{\partial p_k}(p,q)\frac{\partial H}{\partial q_k}(p,q)\right)\\
&=& \sum_k\left(\frac{\partial A}{\partial q_k}(p,q)\dot{q}_k(p,q)
+\frac{\partial A}{\partial p_k}(p,q)\dot{p}_k(p,q)\right)\\
&=&\dot{A}(p,q)\\
\{q_m,p_l\} &=& \sum_k\left(\frac{\partial q_m}{\partial q_k}(p,q)\frac{\partial p_l}{\partial p_k}(p,q)
-\frac{\partial q_m}{\partial p_k}(p,q)\frac{\partial p_l}{\partial q_k}(p,q)\right)\\
&=& \sum_k\delta_{mk}\delta_{lk}\\
&=&\delta_{ml}\\
\{G(p,q),p_l\} &=& \frac{\partial G}{\partial q_l}(p,q)\\
\{G(p,q),q_l\} &=& -\frac{\partial G}{\partial p_l}(p,q)
\end{eqnarray*}
Infinitesimale Transformation der Bewegungsgleichungen:
\begin{eqnarray*}
\frac{d\dot{p}_k}{d\epsilon} &=& 
-\sum_l\frac{\partial^2}{\partial q_k\partial p_l}H\left(p,q\right)\frac{dp_l}{d\epsilon}
-\sum_l\frac{\partial^2}{\partial q_k\partial q_l}H\left(p,q\right)\frac{dq_l}{d\epsilon}\\
\frac{d\dot{q}_k}{d\epsilon} &=& 
+\sum_l\frac{\partial^2}{\partial p_k\partial p_l}H\left(p,q\right)\frac{dp_l}{d\epsilon}
+\sum_l\frac{\partial^2}{\partial p_k\partial q_l}H\left(p,q\right)\frac{dq_l}{d\epsilon}
\end{eqnarray*}
Infinitesimale Transformation:
\begin{eqnarray*}
p^\prime_k &=& p_k+\frac{dp_k}{d\epsilon}\epsilon\\
q^\prime_k &=& q_k+\frac{dq_k}{d\epsilon}\epsilon\\
G^\prime(p^\prime,q^\prime) &=& G\left(p+\frac{dp}{d\epsilon}\epsilon,q+\frac{dq}{d\epsilon}\epsilon\right)\\
&=&G(p,q)+\sum_k\left(
\frac{\partial G}{\partial p_k}\frac{dp_k}{d\epsilon}(p,q)+
\frac{\partial G}{\partial q_k}\frac{dq_k}{d\epsilon}(p,q)
\right)\epsilon\\
&=&G(p,q)+\sum_k\left(
-\{G,q_k\}\frac{dp_k}{d\epsilon}(p,q)+
\{G,p_k\}\frac{dq_k}{d\epsilon}(p,q)
\right)\epsilon
\end{eqnarray*}
Kanonische infinitesimale Transformation:
\begin{eqnarray*}
\{q^\prime_l,p^\prime_m\} &=& \delta_{lm}\\
&=& \left\{q_l+\frac{dq_l}{d\epsilon}\epsilon,p_m+\frac{dp_m}{d\epsilon}\epsilon\right\}\\
&=& \delta_{ml}+\left(\left\{q_l,\frac{dp_m}{d\epsilon}\right\}+\left\{\frac{dq_l}{d\epsilon},p_m\right\}\right)\epsilon\\
&=& \delta_{ml}+\left(-\left\{\frac{dp_m}{d\epsilon},q_l\right\}+\left\{\frac{dq_l}{d\epsilon},p_m\right\}\right)\epsilon\\
&=& \delta_{ml}+\left(\frac{\partial}{\partial p_l}\frac{dp_m}{d\epsilon}+\frac{\partial}{\partial q_m}\frac{dq_l}{d\epsilon}\right)\epsilon\\
\frac{\partial}{\partial p_l}\frac{dp_m}{d\epsilon}&=&-\frac{\partial}{\partial q_m}\frac{dq_l}{d\epsilon}
\end{eqnarray*}
scratch
\begin{eqnarray*}
\frac{\partial}{\partial q_k}L\left(\dot{q}+\frac{d\dot{q}}{d\epsilon}\epsilon,q+\frac{dq}{d\epsilon}\epsilon\right) &=&
\frac{\partial}{\partial q_k}L\left(\dot{q},q\right)
+\sum_l\left(\frac{\partial^2}{\partial q_k\partial \dot{q}_l}L\left(\dot{q},q\right)\frac{d\dot{q}_l}{d\epsilon}\epsilon
+\frac{\partial^2}{\partial q_k\partial q_l}L\left(\dot{q},q\right)\frac{dq_l}{d\epsilon}\epsilon\right)\\
&=&\frac{\partial}{\partial q_k}L\left(\dot{q},q\right)
+\frac{\partial}{\partial q_k}\sum_l\left(\frac{\partial}{\partial \dot{q}_l}L\left(\dot{q},q\right)\frac{d\dot{q}_l}{d\epsilon}\epsilon
+\frac{\partial}{\partial q_l}L\left(\dot{q},q\right)\frac{\partial}{\partial q_k}\frac{dq_l}{d\epsilon}\epsilon\right)\\
&&-\sum_l\left(\frac{\partial}{\partial \dot{q}_l}L\left(\dot{q},q\right)\frac{\partial}{\partial q_k}\frac{d\dot{q}_l}{d\epsilon}\epsilon
+\frac{\partial}{\partial q_l}L\left(\dot{q},q\right)\frac{\partial}{\partial q_k}\frac{dq_l}{d\epsilon}\epsilon\right)\\
&=&\frac{\partial}{\partial q_k}L\left(\dot{q},q\right)
+\frac{\partial}{\partial q_k}\sum_l\left(\frac{\partial}{\partial \dot{q}_l}L\left(\dot{q},q\right)\frac{d\dot{q}_l}{d\epsilon}\epsilon
+\frac{\partial}{\partial q_l}L\left(\dot{q},q\right)\frac{dq_l}{d\epsilon}\epsilon\right)\\
&&-\sum_l\left(\frac{\partial}{\partial \dot{q}_l}L\left(\dot{q},q\right)\frac{\partial}{\partial q_k}\frac{d\dot{q}_l}{d\epsilon}\epsilon
+\frac{d}{dt}\left(\frac{\partial}{\partial \dot{q}_l}L\left(\dot{q},q\right)\right)\frac{\partial}{\partial q_k}\frac{dq_l}{d\epsilon}\epsilon\right)\\
&=&\frac{\partial}{\partial q_k}L\left(\dot{q},q\right)
+\frac{\partial}{\partial q_k}\sum_l\left(\frac{\partial}{\partial \dot{q}_l}L\left(\dot{q},q\right)\frac{d\dot{q}_l}{d\epsilon}\epsilon
+\frac{\partial}{\partial q_l}L\left(\dot{q},q\right)\frac{dq_l}{d\epsilon}\epsilon\right)\\
&&-\frac{d}{dt}\sum_l\left(\frac{\partial}{\partial \dot{q}_l}L\left(\dot{q},q\right)\frac{\partial}{\partial q_k}\frac{dq_l}{d\epsilon}\right)\epsilon
\end{eqnarray*}
\end{subsection}
\end{section}

\end{document}
